% Proposed insertion for Observation_Refinement_Descent_and_Cone_Draft_v0.3.tex
% Recommended location: after the general carrier-splitting factorization and
% before the cover-refinement section.  Bibliography additions are supplied in
% MANUSCRIPT_BIBLIOGRAPHY_ADDITIONS.tex.

\section{All-height carrier splitting under noninteracting exclusive inversions}
\label{sec:exclusive-matching}

The low-height rooted-incidence theorem gives a complete answer when the
weighted height is at most two.  In higher dimensions a universal rooted-forest
criterion is impossible: the same-carrier boundary case already supports
integral torsion and arbitrary finite-poset homotopy types.  Nevertheless, an
all-height exact classification survives when the genuinely incompatible
carrier outcomes do not interact.

Let
\[
A_a=\{s\in Q:(s,a)\in Q^\psi\},\qquad a\in\{0,1\},
\]
and put
\[
P=A_1\setminus A_0,\qquad Z=A_0\setminus A_1,\qquad
C=A_0\cap A_1.
\]
Thus $P$ consists of the $1$-only coarse states, $Z$ of the $0$-only coarse
states, and $C$ of the states that genuinely split.

\begin{definition}[Exclusive inversion graph]
\label{def:exclusive-inversion-graph}
The \emph{exclusive inversion graph} is the bipartite graph
\[
G_\psi=(P\sqcup Z,E_\psi),\qquad
E_\psi=\{(p,z)\in P\times Z:p<z\text{ in }Q\}.
\]
An edge records a coarse comparison whose earlier endpoint admits only the new
outcome $1$ while its later endpoint admits only $0$.
\end{definition}

\subsection{The full-doubling model}

Let
\[
D=Q\times\{0<1\},\qquad K=\Delta D,\qquad L=\Delta Q^\psi.
\]
The actual refined carrier is the induced subposet of $D$ obtained by deleting
exactly
\[
M_0=\{(p,0):p\in P\},\qquad
M_1=\{(z,1):z\in Z\}.
\]
Let $K_P$ be obtained by deleting $M_0$, let $K_Z$ be obtained by deleting
$M_1$, and put
\[
U=K_P\cup K_Z,\qquad L=K_P\cap K_Z.
\]
The complex $K_P$ closes upward to the full top row, while $K_Z$ closes
downward to the full bottom row.  Hence both $K_P\hookrightarrow K$ and
$K_Z\hookrightarrow K$ are homotopy equivalences.  The long exact sequences
of triples together with simplicial excision give
\begin{equation}
H_n(K,L;R)\cong H_{n+1}(K,U;R)
\label{eq:two-colour-shift}
\end{equation}
for every coefficient ring $R$.  Since $D\to Q$ is a homotopy equivalence,
$H_n(K,L;R)$ is the reduced homology of the homotopy cofiber of the refined
projection $\Delta\pi:\Delta Q^\psi\to\Delta Q$.

A simplex of $K\setminus U$ contains at least one missing lower copy $(p,0)$
and at least one missing upper copy $(z,1)$.  Because it is a product-order
chain, every such pair satisfies $p<z$.  Therefore the missing vertices of a
relative simplex project to a complete bipartite subgraph of $G_\psi$.  This
observation isolates the source of higher interaction: bicliques with more than
one edge couple otherwise local defects.

\begin{theorem}[Exclusive-Inversion Matching Theorem]
\label{thm:exclusive-matching}
Suppose that the exclusive inversion graph $G_\psi$ is a matching.  Then, for
every coefficient ring $R$ and every integer $n$,
\begin{equation}
\widetilde H_n\!\left(\Cone(\Delta\pi);R\right)
\cong
\bigoplus_{(p,z)\in E_\psi}
\widetilde H_{n-2}\!\left(
\operatorname{lk}_{\Delta Q}[p,z];R
\right).
\label{eq:exclusive-matching}
\end{equation}
Consequently $\Delta\pi$ is an $R$-homology equivalence if and only if every
coarse edge link $\operatorname{lk}_{\Delta Q}[p,z]$ indexed by an exclusive
inversion is $R$-acyclic.
\end{theorem}

\begin{proof}
A relative simplex of $(K,U)$ contains at least one vertex of $M_0$ and one
vertex of $M_1$.  If it contained two missing lower vertices and one missing
upper vertex, the corresponding upper endpoint would have degree at least two
in $G_\psi$; the dual statement holds for two missing upper vertices.  The
matching hypothesis therefore forces every relative simplex to contain a unique
mixed pair
\[
e=((p,0)<(z,1)),\qquad (p,z)\in E_\psi.
\]
Deleting either endpoint of $e$ sends the simplex into $U$, while deleting any
other vertex preserves the same pair.  Hence the relative chain complex splits
as
\[
C_m(K,U;R)=\bigoplus_{e\in E_\psi}C_m^{(e)},
\qquad
C_m^{(e)}\cong
\widetilde C_{m-2}(\operatorname{lk}_K e;R).
\]
Thus
\[
H_m(K,U;R)\cong
\bigoplus_{e\in E_\psi}
\widetilde H_{m-2}(\operatorname{lk}_K e;R).
\]

It remains to identify the product-order link.  For
$e=((p,0)<(z,1))$, its link factors into the lower interval, the open middle
interval, and the upper interval.  The open middle interval is covered by two
contractible subposets: delete $(z,0)$ and close upward to the bit-$1$ row, or
delete $(p,1)$ and close downward to the bit-$0$ row.  Their intersection is
$(p,z)_Q\times\{0<1\}$, homotopy equivalent to $(p,z)_Q$.  Since $(p,1)$ and
$(z,0)$ are incomparable, every middle-chain lies in one of these two pieces;
therefore
\[
\Delta((p,0),(z,1))\simeq\Sigma\Delta(p,z)_Q.
\]
Joining with the lower and upper factors gives
\[
\operatorname{lk}_{\Delta(Q\times\mathbf2)}e
\simeq
\Sigma\operatorname{lk}_{\Delta Q}[p,z].
\]
Combining this suspension shift with~\eqref{eq:two-colour-shift} yields
\eqref{eq:exclusive-matching}.
\end{proof}

\begin{corollary}[Bounded-poset localization]
\label{cor:bounded-localization}
Assume that $Q$ has a global minimum $\hat0$ and maximum $\hat1$, and that
$G_\psi$ is a matching.  Then every exclusive edge except possibly
$(\hat0,\hat1)$ has contractible coarse edge link.  Hence only the global
endpoint inversion can contribute to the cone homology.
\end{corollary}

\begin{proof}
If $p\ne\hat0$, the lower factor $\Delta Q_{<p}$ has the minimum $\hat0$; if
$z\ne\hat1$, the upper factor $\Delta Q_{>z}$ has the maximum $\hat1$.  In
either case the edge link is a join with a nonempty contractible complex.
\end{proof}

\begin{corollary}[Distributive endpoint profile]
\label{cor:distributive-endpoint}
Let $Q=J(R)$ be the lattice of order ideals of a finite poset $R$.  Make the
bottom state $1$-only, the top state $0$-only, and every other state split.  If
$R$ is an $r$-element antichain, so that $Q=B_r$, then
\[
\widetilde H_r(\Cone(\Delta\pi);\mathbb Z)\cong\mathbb Z
\]
and all other reduced cone homology vanishes.  If $R$ is not an antichain,
the cone is integrally acyclic.
\end{corollary}

\begin{proof}
Corollary~\ref{cor:bounded-localization} reduces the answer to the reduced
homology of the proper part of $J(R)$, shifted by two.  The atom crosscut is a
full simplex unless the atoms join to the top element.  In $J(R)$ the atoms
correspond to the minimal elements of $R$, and they join to the top exactly
when every element of $R$ is minimal, i.e. exactly when $R$ is an antichain.
In that case $J(R)=B_r$ and its proper part is homotopy equivalent to
$S^{r-2}$.
\end{proof}

\subsection{Sharpness: interaction begins at the first shared endpoint}

The matching condition cannot be weakened merely to ``$G_\psi$ is a
forest.''  For the three-element chain $0<1<2$ with availability profile
\[
0:\{1\},\qquad 1:\{0\},\qquad 2:\{0\},
\]
the exclusive graph is the two-edge star $K_{1,2}$.  Direct integral Smith
normal form gives
\[
\widetilde H_1(\Cone(\Delta\pi);\mathbb Z)\cong\mathbb Z,
\]
while the sum of the two individual coarse-edge-link contributions is zero.
The shared endpoint creates a three-vertex interaction core and a differential
between the two edge defects.

Nor does matching imply total unimodularity or torsion-freeness.  Given any
nonempty finite poset $R$, take
\[
Q=\hat0\oplus R\oplus\hat1
\]
and use the endpoint profile above.  There is exactly one exclusive inversion,
and Theorem~\ref{thm:exclusive-matching} gives
\[
\widetilde H_n(\Cone(\Delta\pi);\mathbb Z)
\cong
\widetilde H_{n-2}(\Delta R;\mathbb Z).
\]
Thus arbitrary finite-poset homology, including torsion, can occur even inside
the matching class.  In particular, total unimodularity must be controlled by
stronger conditions on the relevant links; global torsion-freeness alone is
not a TU criterion~\cite{DeyHiraniKrishnamoorthy2011}.

\subsection{Beyond matching: the interaction-core filtration}

For a relative simplex $\sigma\in K\setminus U$, define its missing core
\[
S(\sigma)=\sigma\cap(M_0\cup M_1).
\]
Filter $C_*(K,U)$ by $|S(\sigma)|$.  If $\mathcal C_r$ is the family of mixed
$r$-vertex missing cores that are simplices of $K$ and
\[
\Lambda_S=\operatorname{lk}_K(S)[V(L)],
\]
then the associated graded complex is
\begin{equation}
\operatorname{gr}_r C_m(K,U)
\cong
\bigoplus_{S\in\mathcal C_r}
\widetilde C_{m-r}(\Lambda_S).
\label{eq:interaction-graded}
\end{equation}
Indeed, deleting an available vertex preserves $S$, while deleting a missing
vertex lowers the filtration.  This produces a convergent interaction spectral
sequence whose first page is assembled from the homology of the
$\Lambda_S$.  The cores project to chain-supported bicliques of $G_\psi$.
Under the matching hypothesis only $r=2$ occurs, so the spectral sequence
collapses immediately and recovers Theorem~\ref{thm:exclusive-matching}.

This filtration is the natural next object beyond matching.  The smallest
nonmatching example above shows that shared endpoints already create higher
core differentials; therefore a higher-dimensional analogue of the low-height
rooted-forest theorem must control these interactions rather than only the
acyclicity of $G_\psi$.

\begin{remark}[Priority boundary]
The deletion/cofiber technology used here is classical.  In particular,
vertex-cover obstruction complexes already provide general homotopical tools
for comparing a complex with unions of induced subcomplexes
\cite{ChacholskiEtAl2021}, while poset-fiber decompositions give broad results
under fiber-connectivity hypotheses \cite{BjornerWachsWelker2005}.  The
candidate contribution of Theorem~\ref{thm:exclusive-matching} is the
observation-theoretic two-colour recognition and the resulting exact coarse
edge-link formula.  No claim is made that simplicial vertex deletion or
homotopy cofibers themselves are new.
\end{remark}
